\writeSectionFrame{\presentationTitle}{Fazit}

\realworld{JavaFX}{Szenengraph}
\note{
    Der Szenengraph in JavaFX ist die zentrale Struktur, in der alle UI-Komponenten organisiert werden. Es handelt sich dabei um eine baumartige Struktur, in der jedes Node-Objekt entweder ein Blatt (ein einfaches Element) oder ein Container sein kann, der weitere Node-Objekte enthält. Dies entspricht genau dem Kompositum-Muster, bei dem einzelne Objekte und zusammengesetzte Objekte (Container) gleichermaßen behandelt werden.
}

\realworld{Java DOM API (Document Object Model)}{XML-/HTML-Strukturen}
\note{
    Das DOM (Document Object Model) in Java verwendet das Kompositum-Muster, um die Hierarchie von XML- oder HTML-Dokumenten abzubilden. Elemente wie Knoten (Node), Textinhalte und Attribute werden hierarchisch angeordnet, sodass sowohl einzelne Elemente als auch zusammengesetzte Strukturen gleichermaßen behandelt werden können.
}

\realworld{Java Collections Framework}{Datenstrukturen }
\note{
    Collection ist das Interface, das von konkreten Klassen wie ArrayList oder HashSet implementiert wird. Listen und Mengen können wiederum andere Collections enthalten.
}

\realworld{Apache Commons FileUpload}{Dateiuploads (Multipart-Dateien)}
\note{
    Das Apache Commons FileUpload-Framework verwendet das Kompositum-Muster, um Multipart-Formulareingaben (z.B. Dateiuploads) zu modellieren. Eine hochgeladene Datei kann als ein einzelnes Element betrachtet werden, aber ein Multipart-Formular besteht aus mehreren Dateielementen, die hierarchisch verarbeitet werden.
}

\realworld{Apache POI}{Verarbeitung von Office-Dokumenten (z.B. Excel, Word)}
\note{
    Apache POI verwendet das Kompositum-Muster, um Office-Dokumente wie Excel-Tabellen oder Word-Dokumente zu verarbeiten. Tabellenblätter, Zeilen und Zellen werden als hierarchische Strukturen modelliert.
}

\vorteil{Einfache Abbildung von Hierarchien}
\note{
	Mit dem Kompositum-Muster können Hierarchien sehr einfach abgebildet werden. Auch komplexe Hierarchien sind gut beherrschbar durch die Nutzung von Polymorphie. 
}

\vorteil{Einfache Erweiterbarkeit (Open-Closed-Prinzip)}
\note{
	Neue Komponenten (z.B. neue Typen von Blättern oder Kompositen) können leicht hinzugefügt werden, ohne dass der bestehende Code verändert werden muss. Das fördert die Offenheit für Erweiterungen und die Schließung gegenüber Änderungen (Open-Closed-Prinzip).
}

\vorteil{Flexibilität bei der Strukturierung}
\note{
	Das Kompositum-Muster erlaubt es, komplexe Strukturen aus einfachen Komponenten zu bauen und diese dynamisch zur Laufzeit zu verändern, indem man Komponenten hinzufügt oder entfernt.
}

\vorteil{Reduzierte Komplexität im Clientcode}
\note{
	Da der Client nur mit der gemeinsamen Schnittstelle der Komponenten interagiert, bleibt der Code einfacher und leichter verständlich. Die interne Komplexität der Kompositum-Struktur wird vor dem Client verborgen.
}

\vorteil{Einheitliche Behandlung von Objekten}
\note{
	Das Kompositum-Muster ermöglicht es, einzelne Objekte und zusammengesetzte Objekte (z.B. eine Sammlung von Objekten) auf dieselbe Weise zu behandeln. Dadurch kann der Clientcode vereinfacht werden, da er keine Unterscheidung zwischen Einzel- und Gruppenobjekten treffen muss. Das ist dank Polymorphie möglich. 
}

\nachteil{Komplexität des Designs}
\note{
	Obwohl das Muster die Komplexität im Clientcode reduziert, kann die Implementierung des Musters selbst kompliziert sein. Insbesondere wenn die Hierarchie tief ist oder viele verschiedene Arten von Komponenten existieren, kann das Design schwierig zu handhaben sein.
}

\nachteil{Definition eines allgemeinen Interface kann schwierig sein}
\note{
	Bei Klassen, deren Funktionalitäten sich
	stark unterscheiden, kann die Definition eines allgemeinen Interface schwierig sein. Das kann dazu führen, dass das Interface übergeneralisiert wird.
}

\nachteil{Kann zu inflationärem Umgang verführen}
\note{
	Man sollte das Muster nicht inflationär benutzen. Manchmal sind
	Bäume überhaupt nicht notwendig, sondern flache Listen oder andere einfache Datenstrukturen absolut ausreichend.
}

\nachteil{Schwierigkeiten bei der Durchsetzung von Einschränkungen}
\note{
	Es kann schwierig sein, bestimmte Regeln oder Einschränkungen durchzusetzen, z.B. wenn bestimmte Komponenten nur in bestimmten Kontexten erlaubt sind. Da alle Komponenten dieselbe Schnittstelle haben, kann der Missbrauch dieser Flexibilität Probleme verursachen.
}

\nachteil{Einschränkung der Typensicherheit}
\note{
	Da das Kompositum-Muster eine gemeinsame Schnittstelle für Blätter und Komposita erfordert, kann die Typensicherheit leiden. In stark typisierten Sprachen wie Java kann dies zu weniger aussagekräftigen Typen führen und Fehler möglicherweise erst zur Laufzeit entdeckt werden.
}

\nachteil{Performance-Overhead}
\note{
	In einigen Fällen kann das Durchlaufen und Verarbeiten von großen Kompositum-Strukturen einen Performance-Overhead verursachen, insbesondere wenn viele rekursive Aufrufe involviert sind oder komplexe Operationen durchgeführt werden müssen.
}

\nachteil{Erhöhte Anforderungen an den Speicher}
\note{
	Bei sehr großen Strukturen kann das Muster zu erhöhtem Speicherverbrauch führen, da jedes Objekt eine Referenz auf seine Kinder und potenziell auch auf seinen Elternknoten hält.
}

\begin{frame}{Kompositum einsetzen, wenn ...}
	\begin{itemize}
 		\item eine hierarchische Struktur modelliert werden muss.
 		\item die einheitliche Behandlung von Einzel- und zusammengesetzten Objekten gewünscht ist.
 		\item dynamische Änderungen an der Struktur möglich sein sollen
        \item eine rekursive Verarbeitung der Struktur erforderlich ist
        \item Wiederverwendbarkeit und Erweiterbarkeit im Vordergrund stehen
        \item der Clientcode vereinfacht werden soll, indem die Komplexität der Struktur abstrahiert wird
 	\end{itemize}
\end{frame}

\begin{frame}{Kompositum NICHT einsetzen, wenn ...}
	\begin{itemize}
 		\item keine Hierarchie erforderlich ist
 		\item Einzel- und zusammengesetzte Objekte sehr unterschiedliche Operationen erfordern\textbf{}
 		\item die Struktur einfach und statisch ist
 	\end{itemize}
\end{frame}